\nonstopmode{}
\documentclass[a4paper]{book}
\usepackage[times,inconsolata,hyper]{Rd}
\usepackage{makeidx}
\makeatletter\@ifl@t@r\fmtversion{2018/04/01}{}{\usepackage[utf8]{inputenc}}\makeatother
% \usepackage{graphicx} % @USE GRAPHICX@
\makeindex{}
\begin{document}
\chapter*{}
\begin{center}
{\textbf{\huge Package `GradBoostR'}}
\par\bigskip{\large \today}
\end{center}
\ifthenelse{\boolean{Rd@use@hyper}}{\hypersetup{pdftitle = {GradBoostR: Gradient Boosting + RF + NN Engine}}}{}
\ifthenelse{\boolean{Rd@use@hyper}}{\hypersetup{pdfauthor = {Nic Coxen}}}{}
\begin{description}
\raggedright{}
\item[Type]\AsIs{Package}
\item[Title]\AsIs{Gradient Boosting + RF + NN Engine}
\item[Version]\AsIs{0.1.0}
\item[Description]\AsIs{Fast ML engine with boosting, random forests, and neural nets.}
\item[License]\AsIs{GPL-3}
\item[Encoding]\AsIs{UTF-8}
\item[Imports]\AsIs{Rcpp}
\item[LinkingTo]\AsIs{Rcpp, RcppArmadillo}
\item[Config/roxygen2/version]\AsIs{8.1.0}
\item[NeedsCompilation]\AsIs{yes}
\item[Author]\AsIs{Nic Coxen [aut, cre]}
\item[Maintainer]\AsIs{Nic Coxen }\email{nicjcoxen@gmail.com}\AsIs{}
\item[Archs]\AsIs{x64}
\end{description}
\Rdcontents{Contents}
\HeaderA{gradboostr\_fit}{Train a GradBoostR model}{gradboostr.Rul.fit}
%
\begin{Description}
Train a GradBoostR model
\end{Description}
%
\begin{Usage}
\begin{verbatim}
gradboostr_fit(x, y, method = c("gbm", "rf", "rf_class", "nn"), ...)
\end{verbatim}
\end{Usage}
%
\begin{Arguments}
\begin{ldescription}
\item[\code{x}] Matrix or data.frame of features.

\item[\code{y}] Numeric (regression) or factor/integer (classification).

\item[\code{method}] One of: "gbm", "rf", "rf\_class", "nn".

\item[\code{...}] Additional hyperparameters passed to the underlying C++ engine.
\end{ldescription}
\end{Arguments}
%
\begin{Value}
A gradboostr\_model object.
\end{Value}
\HeaderA{gradboostr\_predict}{Predict using a GradBoostR model}{gradboostr.Rul.predict}
%
\begin{Description}
Predict using a GradBoostR model
\end{Description}
%
\begin{Usage}
\begin{verbatim}
gradboostr_predict(object, newdata, type = c("response", "prob"), ...)
\end{verbatim}
\end{Usage}
%
\begin{Arguments}
\begin{ldescription}
\item[\code{object}] A gradboostr\_model from gradboostr\_fit().

\item[\code{newdata}] Matrix or data.frame of features.

\item[\code{type}] "response" or "prob" (classification only).

\item[\code{...}] Not used.
\end{ldescription}
\end{Arguments}
%
\begin{Value}
Numeric vector (regression), factor (classification), or probability matrix.
\end{Value}
\HeaderA{hello}{Hello, World!}{hello}
%
\begin{Description}
Prints 'Hello, world!'.
\end{Description}
%
\begin{Usage}
\begin{verbatim}
hello()
\end{verbatim}
\end{Usage}
%
\begin{Examples}
\begin{ExampleCode}
hello()
\end{ExampleCode}
\end{Examples}
\printindex{}
\end{document}
